\section{Using the full gap and gradient norm together}
\label{joint:section}
The two observable bounds can be combined more tightly than by taking
their minimum. A bounded gradient norm excludes the shrinking-radius
examples that determine part of the full-gap envelope. The resulting
certificate reuses the same scalar function $V$ on a shorter interval.

\begin{theorem}[The sharp joint certificate]
\label{joint:main}
Assume \eqref{boundarygap:assumptions}. For $g,P>0$ put
\begin{equation}\label{joint:definition}
 h=\min\{g,P^2/m\},\qquad
 \mathcal J(g,P)=
 \max_{h^2/P^2\le\rho\le h/m}V(h,\rho).
\end{equation}
Set $\mathcal J(g,0)=\mathcal J(0,P)=0$. Then
\begin{equation}\label{joint:bound}
 \boxed{\delta(x)\le
 \mathcal J\bigl(G(x),\|\nabla f(x)\|\bigr)
 \le\min\{\mathcal B(G(x)),\mathcal E(\|\nabla f(x)\|)\}.}
\end{equation}
Equivalently,
\[
 \mathcal J(g,P)=
 \begin{cases}
 \displaystyle\max_{g^2/P^2\le\rho\le g/m}V(g,\rho),
       &0<g<P^2/m,\\[2mm]
 \mathcal E(P),&g\ge P^2/m.
 \end{cases}
\]
The maximum exists. The function $\mathcal J$ is finite, continuous
on $[0,\infty)^2$, and nondecreasing in each argument.

For every fixed $g,P>0$, $\mathcal J(g,P)$ is the supremum of the
source error in the stated class with
\[
 G(x)=g,\qquad \|\nabla f(x)\|=P
\]
both held exactly. The same supremum holds with these two equalities
replaced by upper bounds. Smooth near-extremizers exist in dimension
at most two, on compact convex sets with nonempty interior in their
ambient space. The domain is allowed to vary; no fixed-domain or
prescribed-diameter sharpness is asserted.
\end{theorem}

Appendix~\ref{joint:proof} proves the result. If $g\ge P^2/m$, the
interval in \eqref{joint:definition} is a single point. Otherwise,
the lower endpoint $g^2/P^2$ is positive. It removes radii at which
the full-gap extremal trajectory would require a source gradient norm
larger than the observed $P$.

\begin{corollary}[Strict improvement]
\label{joint:strict}
For every $H_1\ge0$ and $0<g<P^2/m$,
\[
 \mathcal J(g,P)\le\mathcal E(P)
       -\frac{(P^2-mg)^2}{2mP^2}<\mathcal E(P).
\]
If additionally $H_1>0$ and $\mathcal S(g)\le0$ in
Theorem~\ref{gapregime:main}, then
\[
 \mathcal J(g,P)<\min\{\mathcal B(g),\mathcal E(P)\}.
\]
\end{corollary}
\begin{proof}
For a radius $\rho\in[g^2/P^2,g/m]$ put
$p=g/\sqrt\rho$ and $q=p-m\sqrt\rho$.
The radial deficit \eqref{bndnorm:radial-deficit}, applied to smooth
near-extremizers, gives $V(g,\rho)+\mathcal E(q)\le\mathcal E(p)$.
The growth inequality in Appendix~\ref{energy:proof} gives
$\mathcal E(P)-\mathcal E(p)\ge(P^2-p^2)/(2m)$ and
$\mathcal E(q)\ge q^2/(2m)$. Hence
\[
 \mathcal E(P)-V(g,\rho)
 \ge\frac{P^2-p^2+q^2}{2m}
 =\frac{P^2-2mg+m^2\rho}{2m}
 \ge\frac{(P^2-mg)^2}{2mP^2}.
\]
Taking the maximum over radii proves the first claim.
In the second case Theorem~\ref{gapregime:main} gives
$V(g,\rho)<g$ at every positive radius. The joint maximum is
over a compact interval with positive lower endpoint, so it is
strictly less than $g=\mathcal B(g)$.
\end{proof}
The quantitative deficit is attained when $H_0=m$ and $H_1=0$.

\begin{corollary}[One root or an elementary joint bound]
\label{joint:ordinary}
Let $H_1>0$, $g,P>0$, and $h=\min\{g,P^2/m\}$.
If $\mathcal S(h)\le0$, then
\begin{equation}\label{joint:ordinary-formula}
 \boxed{\mathcal J(g,P)=V(h,h^2/P^2)
       \le h+\frac{h^2}{P^2}\mathcal S(h)\le h.}
\end{equation}
The exact joint envelope therefore requires only the single
switch root at $\rho=h^2/P^2$. The displayed upper bound uses
only a square root and arithmetic to evaluate $\mathcal S(h)$;
it requires neither a switch root nor an outer grid.
\end{corollary}
\begin{proof}
Corollary~\ref{gapregime:radius-monotone} states that
$V(h,\rho)$ is strictly decreasing in positive $\rho$ and that
$V(h,\rho)\le h+\rho\mathcal S(h)$ in this regime.
Apply both conclusions at the left endpoint of
\eqref{joint:definition}. They also apply when the interval
has collapsed to one point.
\end{proof}

For example, $m=H_0=H_1=1$, $g=1$, and $P=2$ give
\[
 \mathcal S(1)=4-\frac{7\sqrt{14}}6<0,\qquad
 \delta(x)\le\mathcal J(1,2)
 \le2-\frac{7\sqrt{14}}{24}<1.
\]
This is an elementary strengthening of the gap-only bound one.
It follows from the theorem, not from a scan of radii.
The single-root value is also the fixed-initial-slope envelope
in Proposition~\ref{joint:fixed-slope}, with
$q=P-mh/P$ and length $h/P$.
With certified source upper bounds use the same formulas with
$g=B^+$ and $P=P^+$; a numerical shortcut requires a certified
nonpositive sign for $\mathcal S(h)$. An unresolved sign uses
the general valid procedure.

\begin{corollary}[Closed form for constant curvature]
\label{joint:constant}
If $H_1=0$, define
\[
 \beta=\frac{H_0^2-2mH_0-m^2}{4H_0},\qquad
 h=\min\{g,P^2/m\}.
\]
For $g,P>0$,
\begin{equation}\label{joint:constant-formula}
 \mathcal J(g,P)=
 \begin{cases}
 h+\beta h/m,&\beta\ge0,\\
 h+\beta h^2/P^2,&\beta<0.
 \end{cases}
\end{equation}
If $m\le H_0<(1+\sqrt2)m$ and $0<g<P^2/m$, then
\[
 \mathcal J(g,P)<\min\{\mathcal B(g),\mathcal E(P)\}.
\]
For example, $m=H_0=1$, $H_1=0$, $g=1/2$, $P=1$ give
\[
 \mathcal J(1/2,1)=3/8,
 \qquad \mathcal B(1/2)=\mathcal E(1)=1/2.
\]
The quadratic $f(s)=s/2+s^2/2$ on $Q=[0,1/2]$, evaluated at
$s=1/2$, attains the joint value with exactly those observations.
\end{corollary}
\begin{proof}
The constant-curvature expression
$V(h,\rho)=h+\beta\rho$ is affine in $\rho$, so one of the
two interval endpoints gives \eqref{joint:constant-formula}.
In the strict-improvement regime put
$b=-\beta/m\in(0,1/2]$ and $t=mg/P^2\in(0,1)$. Then
\[
 \mathcal B(g)=g,\qquad
 \mathcal J(g,P)=\frac{P^2}{m}(t-bt^2),\qquad
 \mathcal E(P)=\frac{P^2}{m}(1-b).
\]
The first strict comparison follows from $b,t>0$; the second from
$1-b-t+bt^2=(1-t)\{1-b(1+t)\}>0$.
\end{proof}

\paragraph{Certified observations and stopping.}
Suppose a paid source query supplies a full-gap upper bound $B^+$ and
a norm upper bound $P^+$. Monotonicity gives
\[
 \delta(x)\le\mathcal J(B^+,P^+).
\]
For example, a norm-certified gradient with
$\|\widehat g-\nabla f(x)\|\le\eta$ supplies
$P^+=\|\widehat g\|+\eta$. A bound on atom-score differences alone
does not supply this norm bound. The joint certificate uses no
additional objective, gradient, or LMO query when both observations
are already available; its scalar evaluation still has a cost.

Retaining $f(x)-\mathcal J(B^+,P^+)$ strengthens the accumulated
lower bound. Since $\mathcal J(B^+,P^+)\le\mathcal B(B^+)$, all
existing work bounds based on the one-variable calibration
$\mathcal B$ remain valid. There is no need to invert the
two-variable function $\mathcal J$. This is a sharper certificate,
not a claim of fewer optimization steps on untested data.
If the scalar routine returns a certified enclosure
$\mathcal J\le\widehat{\mathcal J}\le\vartheta\mathcal J$,
the same argument uses the effective calibration
$\Phi=\vartheta\mathcal B$ in Theorem~\ref{nloracle:main}.
Its separate oracle-accuracy parameter $\theta$ is unchanged.

\begin{proposition}[Finite scalar evaluation of the joint certificate]
\label{joint:budget}
Suppose $H_1>0$, $g,P>0$, and fix $\vartheta>1$. Let $h$ be as in
\eqref{joint:definition}. Use $K_*=K_*(h)$ and $L_*=L_*(h)$ from
Proposition~\ref{boundarygap:budget}, and put
\[
 \rho_-=h^2/P^2,\qquad \rho_+=h/m,\qquad
 \Delta=\rho_+-\rho_-,\qquad
 n=\left\lceil\log_2\max\left\{1,
                  \frac{8K_*}{m(\vartheta-1)}\right\}\right\rceil.
\]
If $\Delta>0$, set
\[
 J=\max\left\{1,\left\lceil
       \frac{2L_*\Delta}{(\vartheta-1)h}\right\rceil\right\},
 \qquad \rho_j=\rho_-+j\Delta/J,\quad 0\le j\le J.
\]
At each node perform $n$ exact-sign bisections of
\eqref{boundarygap:switch}, with gap argument $h$, and let
$\widehat V_j$ be the upper retained endpoint value. Then
\[
 \widehat{\mathcal J}(g,P)=
 \max_{0\le j\le J}\widehat V_j+\frac{L_*\Delta}{2J}
 \qquad
 \mathcal J(g,P)\le\widehat{\mathcal J}(g,P)
                    \le\vartheta\mathcal J(g,P).
\]
This uses at most $(J+1)n$ root sign comparisons and $J$ maximum
comparisons. If $\Delta=0$, use only its single node, with $n$
root sign comparisons and no grid correction or maximum comparison.
For $H_1=0$, use \eqref{joint:constant-formula}; when either input
is zero return zero. These algebraic branches need no root search.
\end{proposition}
\begin{proof}
The interval contains $h/m$, so
$\mathcal J(g,P)\ge V(h,h/m)=W(h)\ge h/2$.
Each upper endpoint has excess at most
$2\rho_+K_*2^{-n}\le(\vartheta-1)h/4$.
The half-mesh correction has the same upper bound. Their sum is
at most $(\vartheta-1)\mathcal J(g,P)$. The Lipschitz argument of
Proposition~\ref{boundarygap:budget} proves validity, and all
$J+1$ nodes have positive radius. The single-node case uses only
the first error bound.
\end{proof}

These are exact-arithmetic comparison counts. Certified numerical
evaluation must also enclose arithmetic and elementary-function
errors, for instance by reserving a factor
$1<\vartheta_0<\vartheta$. They are not bit-complexity bounds.
Since $h\le g$ and $\Delta\le h/m$, the uniform source-gap cap
\eqref{boundarygap:query-cap} also bounds these budgets: one may use
$J\le\max\{1,\lceil2L_*(\overline B)/(m(\vartheta-1))\rceil\}$
and $n\le n(\overline B)$ at every positive query. Every paid
source query, including unsuccessful refinements and the terminal
check, must be charged for its chosen scalar calculation.
